% Source excerpts: Gessel and Zhuang, arXiv:2306.15785v3 (2024).
% Operator macros and Conjecture 7.1 are reproduced from the source TeX.
\DeclareMathOperator{\pk}{pk}
\DeclareMathOperator{\ipk}{ipk}
\DeclareMathOperator{\lpk}{lpk}
\DeclareMathOperator{\ilpk}{ilpk}
\DeclareMathOperator{\br}{br}
\DeclareMathOperator{\ibr}{ibr}
\DeclareMathOperator{\udr}{udr}
\DeclareMathOperator{\iudr}{iudr}
\DeclareMathOperator{\st}{st}
\DeclareMathOperator{\ist}{ist}

Given a permutation statistic $\st$, let us define its \textit{inverse
statistic} $\ist$ by $\ist(\pi)\coloneqq\st(\pi^{-1})$.

\begin{itemize}
\item The \textit{peak number} $\pk$. We call $i$ (where $2\leq i\leq n-1$)
a \textit{peak} of $\pi\in\mathfrak{S}_{n}$ if $\pi(i-1)<\pi(i)>\pi(i+1)$.
Then $\pk(\pi)$ is the number of peaks of $\pi$.
\item The \textit{left peak number} $\lpk$. We call $i\in[n-1]$ a \textit{left
peak} of $\pi\in\mathfrak{S}_{n}$ if either $i$ is a peak of $\pi$,
or if $i=1$ and $\pi(1)>\pi(2)$. Then $\lpk(\pi)$ is the number
of left peaks of $\pi$.
\item The \textit{number of biruns} $\br$. A \textit{birun} of a permutation
$\pi$ is a maximal monotone consecutive subsequence. Then $\br(\pi)$
is the number of biruns of $\pi$.
\item The \textit{number of up-down runs} $\udr$. An \textit{up-down run}
of $\pi$ is either a birun of $\pi$, or $\pi(1)$ if $\pi(1)>\pi(2)$.
Then $\udr(\pi)$ is the number of up-down runs of $\pi$.\footnote{Equivalently, $\udr(\pi)$ is equal to the length of the longest alternating
subsequence of $\pi$, which was studied in depth by Stanley \cite{Stanley2008}.}
\end{itemize}
\begin{conjecture}
\label{c-rr}The following polynomials are real-rooted for all $n\geq1$:
\begin{enumerate}
\item [\normalfont{(a)}] ${\displaystyle [s^{k+1}]\,P_{n}^{(\pk,\ipk)}(s,t)=\sum_{\substack{\pi\in\mathfrak{S}_{n}\\
\pk(\pi)=k
}
}}t^{\ipk(\pi)+1}$ for  $0\leq k\leq\left\lfloor (n-1)/2\right\rfloor $,
\item [\normalfont{(b)}] ${\displaystyle [s^{k}]\,P_{n}^{(\lpk,\ilpk)}(s,t)=\sum_{\substack{\pi\in\mathfrak{S}_{n}\\
\lpk(\pi)=k
}
}}t^{\ilpk(\pi)}$ for  $0\leq k\leq\left\lfloor n/2\right\rfloor $,
\item [\normalfont{(c)}] ${\displaystyle [s^{k}]\,P_{n}^{(\pk,\ilpk)}(t,s)=\sum_{\substack{\pi\in\mathfrak{S}_{n}\\
\lpk(\pi)=k
}
}}t^{\ipk(\pi)+1}$ for  $0\leq k\leq\left\lfloor n/2\right\rfloor $,
\item [\normalfont{(d)}] ${\displaystyle [s^{k+1}]\,P_{n}^{(\pk,\ilpk)}(s,t)=\sum_{\substack{\pi\in\mathfrak{S}_{n}\\
\pk(\pi)=k
}
}}t^{\ilpk(\pi)}$ for  $0\leq k\leq\left\lfloor (n-1)/2\right\rfloor $,
\item [\normalfont{(e)}] ${\displaystyle [s^{k}]\,P_{n}^{(\udr,\ipk)}(s,t)=\sum_{\substack{\pi\in\mathfrak{S}_{n}\\
\udr(\pi)=k
}
}}t^{\ipk(\pi)+1}$ for  $1\leq k\leq n$,
\item [\normalfont{(f)}] ${\displaystyle [s^{k}]\,P_{n}^{(\udr,\ilpk)}(s,t)=\sum_{\substack{\pi\in\mathfrak{S}_{n}\\
\udr(\pi)=k
}
}}t^{\ilpk(\pi)}$ for  $1\leq k\leq n$, and
\item [\normalfont{(g)}] ${\displaystyle [s^{k}]\,P_{n}^{(\br,\ipk)}(s,t)=\sum_{\substack{\pi\in\mathfrak{S}_{n}\\
\br(\pi)=k
}
}}t^{\ipk(\pi)+1}$ for  $1\leq k\leq n$.
\end{enumerate}
\end{conjecture}
